\documentclass[10pt,a4paper]{article}

\usepackage[T1]{fontenc}
\usepackage[utf8]{inputenc}
\usepackage[ngerman]{babel}
\usepackage[a4paper,margin=14mm]{geometry}
\usepackage{lmodern}
\usepackage[table]{xcolor}
\usepackage{tabularx}
\usepackage{amsmath}
\usepackage{microtype}
\usepackage{tikz}
\usetikzlibrary{arrows.meta,calc}

% Die Lösungsfassung unterscheidet sich nur durch \solutiontrue.
\newif\ifsolution
\solutionfalse

\pagestyle{empty}
\setlength{\parindent}{0pt}
\setlength{\parskip}{0pt}
\renewcommand{\familydefault}{\sfdefault}

\definecolor{Navy}{HTML}{17324D}
\definecolor{Blue}{HTML}{3B6EA5}
\definecolor{Sky}{HTML}{EAF3FA}
\definecolor{Gold}{HTML}{F3B33D}
\definecolor{Ink}{HTML}{1E2935}
\definecolor{Muted}{HTML}{617080}
\definecolor{Line}{HTML}{D8E1E8}
\color{Ink}

\newcommand{\FG}{F_{\mathrm G}}
\newcommand{\kopf}[1]{%
  \colorbox{Navy}{\parbox{\dimexpr\linewidth-2\fboxsep\relax}{%
    \vspace{1.3mm}\color{white}
    \begin{tabularx}{\linewidth}{@{}Xr@{}}
      {\small\bfseries PHYSIK · KLASSE 11} &
      \colorbox{Gold}{\color{Navy}\small\bfseries\strut\ifsolution LÖSUNGEN\else ABFRAGE\fi}
    \end{tabularx}
    \vspace{0.9mm}
    {\fontsize{16}{19}\selectfont\bfseries Aufgabe 8 -- Das Gravitationsgesetz}\par
    \vspace{0.3mm}{\small #1}\vspace{1.1mm}
  }}\par\vspace{1.5mm}}
\newcommand{\frage}[2]{%
  \vspace{1.2mm}{\color{Blue}\bfseries\large #1 · #2}\par
  \vspace{0.4mm}{\color{Blue}\rule{\linewidth}{0.65pt}}\par\vspace{0.8mm}}
\newcommand{\antwortfeld}[2]{%
  \par\vspace{0.7mm}%
  \pgfmathsetmacro{\aw}{\linewidth/1pt}%
  \pgfmathsetmacro{\tw}{\aw-7}%
  \begin{tikzpicture}[x=1pt,y=1pt]
    \path[use as bounding box] (0,0) rectangle (\aw,-25.607*#1);
    \ifsolution
      \node[anchor=north west,text width=\tw pt,inner sep=0pt,align=left,
        font=\small,text=Blue] at (0,-0.5) {#2};
    \else
      \foreach \i in {1,...,#1}{\draw[Line,line width=.55pt] (0,-25.607*\i) -- (\aw,-25.607*\i);}
    \fi
  \end{tikzpicture}\par}
\newcommand{\fuss}[1]{%
  \vfill{\color{Line}\rule{\linewidth}{0.5pt}}\par\vspace{1mm}
  \begin{tabularx}{\linewidth}{@{}Xr@{}}
    {\small\color{Muted}Klasse 11 · Physik · Kreisbewegungen} &
    {\small\color{Muted}\ifsolution Lösungen\else Abfrage\fi{} · Seite #1 von 2}
  \end{tabularx}}

\begin{document}
\kopf{Gravitation · Kräfte zwischen zwei Körpern}

\frage{1}{Newtons Gesetz}
\textbf{Nenne} das Gesetz, welches Newton formulierte. \textbf{Erkläre} dessen Bedeutung.
\antwortfeld{2}{Newton formulierte das allgemeine Gravitationsgesetz: Alle Körper mit Masse ziehen einander an. Es gilt auf der Erde und im Weltall; beispielsweise ist die Anziehung der Erde auf den Mond die Kraft, die ihn auf seiner Bahn umlenkt.}

\frage{2}{Zwei Körper}
\begin{center}
\begin{tikzpicture}[x=1mm,y=1mm]
  \path[use as bounding box] (24,-3) rectangle (147,38);
  \coordinate (A) at (37,18);
  \coordinate (B) at (133,18);
  \draw[Navy,line width=.9pt,fill=Sky] (A) circle (8);
  \draw[Navy,line width=.9pt,fill=Sky] (B) circle (12);
  \fill[Navy] (A) circle (1.1);
  \fill[Navy] (B) circle (1.1);
  \node[font=\small\bfseries,text=Navy] at (37,24) {A};
  \node[font=\small\bfseries,text=Navy] at (133,28) {B};
  \node[font=\small] at (37,1) {$m_A$};
  \node[font=\small] at (133,1) {$m_B>m_A$};
  \draw[Muted,dashed,line width=.6pt] (37,18) -- (133,18);
  \ifsolution
    \draw[-{Latex[length=2.4mm]},Blue,line width=1.1pt] (37,18) -- (62,18);
    \draw[-{Latex[length=2.4mm]},red!80!black,line width=1.1pt] (133,18) -- (108,18);
    \node[fill=white,inner sep=1pt,font=\scriptsize,text=Blue] at (49,35) {$F_{\mathrm{B\ auf\ A}}$};
    \node[fill=white,inner sep=1pt,font=\scriptsize,text=red!80!black] at (121,35) {$F_{\mathrm{A\ auf\ B}}$};
  \fi
\end{tikzpicture}
\end{center}
\textbf{Zeichne} an beiden Körpern die gegenseitigen Gravitationskräfte als beschriftete Pfeile ein. \textbf{Vergleiche} ihre Beträge und Richtungen. \textbf{Beziehe Stellung}, ob, und falls ja, welche Körper beschleunigt werden.
\antwortfeld{3}{Die beiden Kräfte sind gleich groß und entgegengesetzt gerichtet, wirken aber auf verschiedene Körper: $F_{\mathrm{A\ auf\ B}}$ wirkt auf B, $F_{\mathrm{B\ auf\ A}}$ auf A. Beide Körper werden zum jeweils anderen Körper hin beschleunigt, da auf jeden eine Gravitationskraft wirkt. Bei gleichem Kraftbetrag wird der leichtere Körper A nach $a=\frac{F}{m}$ stärker beschleunigt als B.}

\frage{3}{Änderungen}
Ausgehend von unveränderten Anfangswerten wird jede Zeile \textbf{unabhängig} betrachtet; die jeweils nicht genannten Größen bleiben gleich. \textbf{Trage} den Faktor des neuen Kraftbetrags ein und \textbf{begründe} den Fall des doppelten Abstands.
\par\vspace{0.8mm}
\renewcommand{\arraystretch}{1.48}
\begin{tabularx}{\linewidth}{|>{\raggedright\arraybackslash}X|>{\centering\arraybackslash}p{58mm}|}
\hline
\rowcolor{Sky}\textbf{Änderung} & \textbf{Neuer Kraftbetrag im Vergleich zu $F$} \\
\hline
Nur Masse von A verdoppeln & \rule{0pt}{9mm}\ifsolution $2F$\fi \\
\hline
Nur Masse von B verdoppeln & \rule{0pt}{9mm}\ifsolution $2F$\fi \\
\hline
Beide Massen verdoppeln & \rule{0pt}{9mm}\ifsolution $4F$\fi \\
\hline
Mittelpunktabstand verdoppeln & \rule{0pt}{9mm}\ifsolution $\frac{F}{4}$\fi \\
\hline
\end{tabularx}
\ifsolution
\antwortfeld{2}{Bei doppeltem Mittelpunktabstand steht $(2r)^2=4r^2$ im Nenner. Der neue Kraftbetrag beträgt daher ein Viertel des ursprünglichen Betrags.}
\else
\antwortfeld{2}{}
\fi

\fuss{1}
\newpage

\kopf{Gravitation · Kraft berechnen}

\frage{4}{Kraft berechnen}
Zwei kugelsymmetrische Körper haben die Massen $m_1=20{,}0\,\mathrm{kg}$ und $m_2=30{,}0\,\mathrm{kg}$; ihre Mittelpunkte liegen $r=2{,}00\,\mathrm{m}$ auseinander. \textbf{Berechne} den gegenseitigen Kraftbetrag mit Rechenweg und Einheit; verwende $G=6{,}67430\cdot10^{-11}\,\frac{\mathrm N\cdot\mathrm m^2}{\mathrm{kg}^2}$ und beachte die Anzahl gültiger Ziffern. Du darfst die Formel $\FG=G\,\frac{m_1m_2}{r^2}$ verwenden.
\antwortfeld{8}{Ansatz: $\displaystyle \FG=G\,\frac{m_1m_2}{r^2}$.\\[1mm]
Einsetzen: $\displaystyle \FG=\frac{6{,}67430\cdot10^{-11}\,\frac{\mathrm N\cdot\mathrm m^2}{\mathrm{kg}^2}\cdot20{,}0\,\mathrm{kg}\cdot30{,}0\,\mathrm{kg}}{(2{,}00\,\mathrm m)^2}=1{,}001145\cdot10^{-8}\,\mathrm N$.\\[1mm]
Die Messwerte $20{,}0$, $30{,}0$ und $2{,}00$ haben je drei gültige Ziffern; G ist genauer angegeben. Erst am Ende gerundet: $\boxed{\FG=1{,}00\cdot10^{-8}\,\mathrm N}$.}

\fuss{2}
\end{document}
